\chapter*{Abstract}
\addcontentsline{toc}{chapter}{Abstract}
\markboth{Abstract}{Abstract}
\begin{spacing}{1.0}
\begin{otherlanguage}{english}

The scarcity of labelled electrocardiograms (ECGs) for rare conditions makes it difficult to train predictive models, which can lead to clinically important problems, such as Brugada syndrome, being left unaddressed by machine learning. This undergraduate thesis investigates whether Gemma~4, an open-weight vision-language model, can classify ECG images to identify possible cases of Brugada syndrome through in-context learning. The model receives instructions and several labelled examples in the prompt while keeping its weights fixed. A supervised ResNet18 convolutional neural network serves as a baseline for assessing how each method uses the available cases.

The evaluation uses two datasets. The first comprises 100 synthetic cases generated by a simulator that reproduces heartbeats with different combinations of three morphological features. The second contains clinical records provided by the Central University Hospital of Asturias (HUCA) for 317 patients. For each clinical record, a beat is selected from the central portion of the V1 precordial lead recording, with the patient’s clinical label used as the reference. The two datasets make it possible to examine how the source of the examples affects classification.

Both Gemma~4 and ResNet18 improve as they receive more examples from the simulator, but they struggle to generalise to clinical ECGs from HUCA, even with the domain adaptation strategies evaluated. With demonstrations from HUCA, Gemma~4 outperforms ResNet18 with eight or more examples and reaches a balanced accuracy of 0.761 with 32 cases split equally between the two classes. In the evaluation without demonstrations, Gemma~4 did not identify any positive cases. The improvement occurs when the examples match the task presented by the query ECG.

This result motivates further research into in-context learning to prioritise the review of suspicious ECGs and facilitate data collection for developing models with better performance or explainability. The use of an open-weight model also allows inference to run locally without sending clinical data to external services, helping to protect patient privacy. Clinical validation would be needed to establish the usefulness of this approach in supporting medical assessment.

\textbf{Keywords:} ECG images, Brugada syndrome, in-context learning, vision-language models, convolutional neural networks, domain adaptation.

\end{otherlanguage}
\end{spacing}
